% quiver style
\usepackage{tikz-cd}
% `calc` is necessary to draw curved arrows.
\usetikzlibrary{calc}
% `pathmorphing` is necessary to draw squiggly arrows.
\usetikzlibrary{decorations.pathmorphing}

% A TikZ style for curved arrows of a fixed height, due to AndréC.
\tikzset{curve/.style={settings={#1},to path={(\tikztostart)
					.. controls ($(\tikztostart)!\pv{pos}!(\tikztotarget)!\pv{height}!270:(\tikztotarget)$)
					and ($(\tikztostart)!1-\pv{pos}!(\tikztotarget)!\pv{height}!270:(\tikztotarget)$)
					.. (\tikztotarget)\tikztonodes}},
	settings/.code={\tikzset{quiver/.cd,#1}
			\def\pv##1{\pgfkeysvalueof{/tikz/quiver/##1}}},
	quiver/.cd,pos/.initial=0.35,height/.initial=0}

% TikZ arrowhead/tail styles.
\tikzset{tail reversed/.code={\pgfsetarrowsstart{tikzcd to}}}
\tikzset{2tail/.code={\pgfsetarrowsstart{Implies[reversed]}}}
\tikzset{2tail reversed/.code={\pgfsetarrowsstart{Implies}}}
% TikZ arrow styles.
\tikzset{no body/.style={/tikz/dash pattern=on 0 off 1mm}}

% useful macro for class
\newcommand{\at}[3]{\left.#1\right\vert_{#2}^{#3}}
\newcommand\quotient[2]{
	\mathchoice
	{% \displaystyle
		\text{\raise1ex\hbox{$#1$}\Big/\lower1ex\hbox{$#2$}}%
	}
	{% \textstyle
		#1\,/\,#2
	}
	{% \scriptstyle
		#1\,/\,#2
	}
	{% \scriptscriptstyle
		#1\,/\,#2
	}
}

\newcommand{\True}{\textsf{True}}
\newcommand{\T}{\textsf{T}}
\newcommand{\False}{\textsf{False}}
\newcommand{\F}{\textsf{F}}
\newcommand{\act}{\rotatebox[origin=c]{-180}{\(\,\circlearrowright\,\)}}

\usepackage{physics}
\usepackage{complexity}

\DeclareMathOperator{\id}{id}
\DeclareMathOperator{\Span}{span}
\DeclareMathOperator{\supp}{supp}
\DeclareMathOperator{\vol}{vol}
\DeclareMathOperator{\dist}{dist}
\DeclareMathOperator{\diam}{diam}
\DeclareMathOperator{\im}{Im}
\DeclareMathOperator{\sgn}{sgn}
\DeclareMathOperator{\Int}{Int}
\DeclareMathOperator{\diag}{diag}
\DeclareMathOperator{\dom}{dom}
\DeclareMathOperator*{\argmax}{arg\,max}
\DeclareMathOperator*{\argmin}{arg\,min}
\DeclareMathOperator{\Var}{Var}
\DeclareMathOperator{\Cov}{Cov}
\DeclareMathOperator{\Corr}{Corr}

\let\implies\Rightarrow
\let\impliedby\Leftarrow
\let\iff\Leftrightarrow

\newcommand*{\vv}[1]{\vec{\mkern0mu#1}}

\usepackage{stmaryrd} % for \lightning
\newcommand\conta{\scalebox{1.1}{\(\lightning\)}}

\usepackage{bm}
\usepackage{bbm}

% figure support
\usepackage{import}
\usepackage{xifthen}
\ifdefined\pdfminorversion
  \pdfminorversion=7
\fi

\usepackage{pdfpages}
\usepackage{transparent}
\newcommand{\incfig}[2][\columnwidth]{%
	\def\svgwidth{#1}
	\import{Figures/}{#2.pdf_tex}
}

\makeatletter
\def\moverlay{\mathpalette\mov@rlay}
\def\mov@rlay#1#2{\leavevmode\vtop{%
   \baselineskip\z@skip \lineskiplimit-\maxdimen
   \ialign{\hfil$\m@th#1##$\hfil\cr#2\crcr}}}
\newcommand{\charfusion}[3][\mathord]{
    #1{\ifx#1\mathop\vphantom{#2}\fi
        \mathpalette\mov@rlay{#2\cr#3}
      }
    \ifx#1\mathop\expandafter\displaylimits\fi}
\makeatother

\newcommand{\cupdot}{\charfusion[\mathbin]{\cup}{\cdot}}
\newcommand{\bigcupdot}{\charfusion[\mathop]{\bigcup}{\cdot}}

\definecolor{match01}{RGB}{228,26,28}
\definecolor{match02}{RGB}{55,126,184}
\definecolor{match03}{RGB}{77,175,74}
\definecolor{match04}{RGB}{152,78,163}
\definecolor{match05}{RGB}{255,127,0}
\definecolor{match06}{RGB}{255,255,51}
\definecolor{match07}{RGB}{166,86,40}
\definecolor{match08}{RGB}{247,129,191}
\definecolor{match09}{RGB}{153,153,153}
\definecolor{match10}{RGB}{27,158,119}
\definecolor{match11}{RGB}{217,95,2}
\definecolor{match12}{RGB}{117,112,179}
\definecolor{match13}{RGB}{231,41,138}
\definecolor{match14}{RGB}{102,166,30}
\definecolor{match15}{RGB}{230,171,2}
\definecolor{match16}{RGB}{166,118,29}
\definecolor{match17}{RGB}{102,102,102}
\definecolor{match18}{RGB}{141,211,199}
\definecolor{match19}{RGB}{255,255,179}
\definecolor{match20}{RGB}{190,186,218}
\definecolor{match21}{RGB}{251,128,114}

\newcommand{\matchingcolor}[1]{match\ifnum#1<10 0\fi\number#1}

\newcommand{\drawmatching}[1]{%
    \begin{scope}
        \def\R{0.95}
        \coordinate (v21) at (0,0);
        \foreach \i in {0,...,20}{
            \coordinate (v\i) at ({90-360*\i/21}:\R);
        }

        \node[font=\scriptsize] at (0,1.42) {Week \number\numexpr #1+1\relax};

        \draw[gray!45, thin] (0,0) circle (\R);
        \foreach \i in {0,...,20}{
            \fill[black] (v\i) circle (0.028);
        }
        \fill[black] (v21) circle (0.028);

        \draw[thick, RoyalBlue] (v21) -- (v#1);
        \foreach \k in {1,...,10}{
            \pgfmathtruncatemacro{\a}{mod(#1+\k,21)}
            \pgfmathtruncatemacro{\b}{mod(#1-\k+21,21)}
            \draw[thick, RoyalBlue] (v\a) -- (v\b);
        }
    \end{scope}%
}

\newcommand{\drawallmatchings}{%
    \def\R{4.15}
    \coordinate (v21) at (0,0);
    \foreach \i in {0,...,20}{
        \coordinate (v\i) at ({90-360*\i/21}:\R);
    }

    \foreach \r in {0,...,20}{
        \pgfmathtruncatemacro{\week}{\r+1}
        \draw[line width=0.8pt, draw=\matchingcolor{\week}] (v21) -- (v\r);
        \foreach \k in {1,...,10}{
            \pgfmathtruncatemacro{\a}{mod(\r+\k,21)}
            \pgfmathtruncatemacro{\b}{mod(\r-\k+21,21)}
            \draw[line width=0.8pt, draw=\matchingcolor{\week}] (v\a) -- (v\b);
        }
    }

    \draw[black!35, thin] (0,0) circle (\R);
    \foreach \i in {0,...,20}{
        \fill[black] (v\i) circle (0.055);
    }
    \fill[black] (v21) circle (0.065);
}